\begin{exercise}[Electric weights, magnetic cocharacters, and the global gauge group]{I-CH10-020}
Let $G_\Gamma=\bar G/\Gamma$, where $\bar G$ is compact and simply connected
and $\Gamma\subset Z(\bar G)$.  Suppose Higgsing leaves an unbroken group
$H$ with maximal torus $T$.  Define the electric character lattice
$X^*(T)=\operatorname{Hom}(T,U(1))$ and the magnetic cocharacter lattice
$X_*(T)=\operatorname{Hom}(U(1),T)$.  Show that their natural pairing is
integral and reproduces the Dirac condition for every allowed electric weight.
For dyons, derive the antisymmetric pairing that generalizes
$e_i g_j-e_j g_i=2\pi n_{ij}$.

Track the global form of the gauge group.  Show that the allowed weights of
$G_\Gamma$ are precisely the weights of $\bar G$ on which every element of
$\Gamma$ acts trivially.  Describe a magnetic cocharacter as a path whose lift
to $\bar G$ ends in $\Gamma$.  Combine this with
\[
  \pi_2(G_\Gamma/H)\simeq
  \ker\!\left[\pi_1(H)\to\pi_1(G_\Gamma)\right]
\]
to distinguish all Dirac-allowed magnetic charges from those realized by
nonsingular monopoles.  Work out the $SU(2)$ versus $SO(3)=SU(2)/\mathbb Z_2$
example, including fundamental and adjoint electric weights and the even-loop
restriction in the quotient description.
\end{exercise}
